\section{Adding synchronous nonlinear scans}
\label{sec:scans}

The previous argument concerns permutations and pointwise Boolean gates.
We now allow information to travel along a bank through a carried state.
The resulting obstruction uses one global record of exceptional bits
and scan boundary states. It does not require the intermediate banks to
remain independent or linearly encoded.

\begin{definition}[Public scan program]
\label{def:scan-program}
A program has $q$ banks of $P$ bits and $L$ layers. Each initial bank is
either a copy of the raw input or a fixed public bit string. Each layer
first applies one permitted stage from \cref{def:geometry-stage} to each
bank, then scans all banks synchronously from left to right or from
right to left. At an address the scan reads the $q$ old bits and its
current state, writes $q$ new bits, and updates the state. The transition
rule may be any public address-dependent nonlinear function.

The state is carried between layers. Its finite alphabet has at most
$2^s$ values, for an integer $s\geq0$, and its initial value is fixed.
The stages, directions, transition tables and initial constants are
independent of the input. The output is one designated final bank.
\end{definition}

\lean{RawScanProgram} gives these raw-address semantics. Grouping a scan
into analytical bins preserves every output bit and its terminal state
by \lean{RawScanProgram.run_onGrid}. In particular, the program is fixed
before the proof selects a bin scale. Public choices may vary with $d$;
the theorem below fixes only $q,L,s$.

We first isolate a counting argument. The input ranges over the full
Boolean cube, with its coordinates partitioned into public source bins.
A \emph{named demand} is an original input bit $x_v$, together with a
public set $A_v$ of source bins. A helper $H(x)$ determines that demand
if equal helper values and equal input bits in $A_v$ force equal values
of $x_v$, for every pair of inputs.

\begin{lemma}[One shared helper]
\label{lem:named-helper}
Let $F$ be a set of distinct original input coordinates, each with a
named demand determined by one common helper and its own set $A_v$. Suppose each demand's source bin is absent from
$A_v$ and $|A_v|\leq D$ for $v\in F$. If the helper has at most $2^h$
possible values, where $h\in\N$, then
\begin{equation}
 \label{eq:named-helper}
                       |F|\leq (D+1)h.
\end{equation}
Only demands in $F$ need to be determined.
\end{lemma}

\begin{proof}
Order the source bins uniformly at random. Select $v\in F$ when its
source bin precedes every bin in $A_v$. Its selection probability is
$1/(|A_v|+1)\geq1/(D+1)$.

Fix an order and restrict to inputs whose nonzero coordinates are all
selected demands. Two such inputs with equal helpers agree on every
selected bit: work backwards through the bin order. A selected demand's
ancestor bins come later, so their selected bits have already been
recovered; their unselected bits are zero. Local determination now
recovers that demand. No selected demand uses its own bin as local side information.
Thus the helper is injective on this subcube. If it has $t$ selected
coordinates, then $2^t\leq2^h$, giving $t\leq h$. Averaging yields
$|F|/(D+1)\leq h$.
\end{proof}

This is the acyclic-side-information method from index coding
\cite[Theorem~7]{BBJK2011}, specialized to public bins. We do not claim a new
information principle. The Lean result \lean{helper_named_demands_bound}
uses one range-cardinality budget, without assuming linearity or
independence between helper fields. It formalizes the averaging by an
exact finite count of bin orders, followed by the descending-rank
injection above. The budget $h$ is an integer; a real entropy budget is
not asserted. The scan-specific work is to construct this common helper
and prove the required public ancestry bounds.

For a scale $B\mid d$, put $n=P/B$ and record the current bit at every
bad source position before each layer's permutations, in every bank. Let their total number be
$E_B$. Also record the actual state entering each of the $n$ bins in every
scan, in its scan direction. This includes the inherited initial state.
A scan's terminal state is not needed to replay its outputs; if another
layer follows, that state is recorded as the next scan's initial state. The complete
record has at most
\begin{equation}
 \label{eq:scan-record}
                    2^{E_B+sLn}
\end{equation}
possible values. This is a cardinality upper bound; the recorded fields
can be dependent.

Given this record, a final output bin depends on at most
$D=(11q)^L$ original bins. The induction uses at most $11$ good predecessor bins per
permuted bank and the scan's entering state. A scan's output on one bin
is determined by that state and the old tuples in the bin, in either
direction. Both local replay and the entire layered reconstruction are
proved from the transition recurrence in
\lean{Scan.output_congr_of_record} and
\lean{ScanProgram.run_congr_on_ancestry}. The ancestor sets are public
and independent of the record.

\begin{theorem}[Finite scan bound]
\label{thm:scan-bound}
Suppose a public scan program transposes every $d\times d$ bit matrix,
where $d\geq1$.
For any $m\geq1$ positive divisor scales $B_i\mid d$ satisfying
$V(B_i)\leq B_j$ for $i<j$, there is a common scale $B=B_i$ such that
\begin{equation}
 \label{eq:scan-bound}
 \frac{P-DP/B}{D+1}
 \leq \frac{qLP}{m}+\frac{sLP}{B},
 \qquad P=d^2,\quad D=(11q)^L.
\end{equation}
\end{theorem}

\begin{proof}[Proof outline]
Each output bin requests $B$ bits from distinct source bins under
transpose. At most $D$ of them come from its own ancestor bins, leaving
at least $P-DP/B$ far demands overall. Apply \cref{lem:named-helper}
with the complete record in \eqref{eq:scan-record}. This gives the
one-scale inequality with $E_B$ in place of $qLP/m$.

The one-stage volume bound, summed over all $q$ banks and $L$ layers,
gives $\sum_i E_{B_i}\leq qLP$. Choose one scale attaining the average.
No volume is pulled back through nonlinear scans. The compiled endpoint
is \lean{RawScanProgram.exists_scale_transpose_bound}.
\end{proof}

\begin{corollary}[Bounded-state scan exclusion]
\label{cor:scan-exclusion}
For fixed $q,L,s$, there is a positive dyadic integer $d_0$ such that
no public scan program with those resource bounds transposes every
$d\times d$ bit matrix at any positive dimension divisible by $d_0$.
In particular, all sufficiently large dyadic dimensions are excluded.
The layouts and transition tables may vary with the dimension.
\end{corollary}

\begin{proof}[Proof outline]
Choose a positive number of scales with
$m\geq 8(D+1)qL$. Choose a dyadic first scale satisfying
\[
 B_0\geq\max\{2D,\,8(D+1)sL,\,2\},
 \qquad B_{i+1}=V(B_i).
\]
Take $d_0=B_m$. Every selected scale divides $d_0$, hence also
every positive multiple $d$ of $d_0$. The left side
of \eqref{eq:scan-bound} is at least $P/[2(D+1)]$. Each term
on the right is at most $P/[8(D+1)]$, so their sum is at most
$P/[4(D+1)]$, a contradiction.
\lean{RawScanProgram.not_transposes_of_threshold_dvd} proves the
positive-multiple exclusion, and
\lean{RawScanProgram.not_transposes_large_dyadic} derives the dyadic
case. \lean{RawScanProgram.not_transposes_of_scales} states the
finite conditions directly.
\end{proof}

The carried state and its boundary records are charged explicitly.
The theorem permits arbitrary nonlinear transition tables, but still
requires fixed public geometry, a bounded number of banks and layers,
and a bounded state alphabet. It does not describe a general multitape
machine with independently moving heads or growing auxiliary workspace.
A construction using those freedoms is not excluded by this result.
